% !PRINTED 536
\weberpage{536}
%<<CONT
die gleichfalls die Diskriminante $-4m$ hat, und worin, wenn $m \equiv 1$
oder $\equiv 2\ (\mathrm{mod}\ 4)$ ist, die beiden äußeren Koeffizienten
ungerade sind; denn es ist:
\[
\begin{aligned}
&\text{für } m \equiv 1\ (\mathrm{mod}\ 4), \quad B \equiv 0\ (\mathrm{mod}\ 2), \quad A+C \equiv 2\ (\mathrm{mod}\ 4),\\
&\text{für } m \equiv 2\ (\mathrm{mod}\ 4), \quad B \equiv 1\ (\mathrm{mod}\ 2), \quad A+C \equiv 0\ (\mathrm{mod}\ 4).
\end{aligned}
\]

Daraus folgt, daß
\[
f(\omega) \quad\text{und}\quad f\left(\dfrac{\omega-1}{\omega+1}\right),
\]
von 24sten Einheitswurzeln abgesehen, dieselbe Wertreihe durchläuft.
Denn ersetzt man $\omega'$ durch eine äquivalente Zahl, so muß, wenn die
äußeren Koeffizienten ungerade bleiben sollen, die Substitution zur
ersten oder zweiten Klasse (§ 36) gehören, und daraus folgt aus (4), daß
auch $\omega$ in eine äquivalente Zahl übergeht. Wenn aber zwei Formen
(5) äquivalent sind, so sind auch die entsprechenden Formen (3) äquivalent
und umgekehrt.

Demnach ist
\[
\Pi\, f(\omega)\, f\!\left(\frac{\omega-1}{\omega+1}\right) = 2^{2h\tau},
\]
andererseits ist aber [§ 34, (18)]:
\[
f(\omega)\, f\!\left(\frac{\omega-1}{\omega+1}\right) = \sqrt{2},
\]
woraus sich ergibt:
\weq{6}{\tau = \frac{1}{4}, \quad m \equiv 1,\, 2\ (\mathrm{mod}\ 4),}
in Übereinstimmung mit dem Resultat des vorigen Paragraphen [§ 141, (31)].

2. Ist sodann $m \equiv 3\ (\mathrm{mod}\ 8)$, so entsprechen einer Wurzel
$\omega'$ einer Form der Diskriminante $-m$ je drei Wurzeln von Formen der
Diskriminante $-4m$:
\[
2\omega', \quad \frac{\omega'}{2}, \quad \frac{\omega'+1}{2},
\]
und es sind die 24sten Potenzen der Größen:
\[
\begin{aligned}
f_1(2\omega') &= \frac{\sqrt{2}}{f_2(\omega')},\\
f_2\!\left(\frac{\omega'}{2}\right) &= \frac{\sqrt{2}}{f_1(\omega')},\\
f_2\!\left(\frac{\omega'+1}{2}\right) &= \frac{\sqrt{2}}{f(\omega')},
\end{aligned}
\]
%CONT>>
